\section{Conclusion}\label{sec:conclusion}

\Argorix demonstrates how agent declarations, policy, sovereign metadata,
provider boundaries, and evidence requirements can share a compiled systems
boundary.  In an authorized observational snapshot, 27 of 33 request
directories were complete, six were source-only, and all 27 complete
EvidenceBundles verified bytecode, trace, security-report, and
trace-event-ledger digest consistency offline; source was outside verification.
The same complete sessions contained 1,188
detailed policy violations and uniformly required review despite a favorable
top-level status.  That contradiction is not incidental: it establishes the
paper's central distinction between integrity and approval.

The implemented contribution is a validated language, fail-closed VM path,
simulated-only executable-provider boundary, runtime trace, SecurityReport,
EvidenceBundle, and semantic digest verification.  Passports and
\code{ans_name} are implemented as declarative identifiers and metadata.
ATrust, MCP, and A2A relationships are validated declarations, not live trust
or protocol execution.  Sovereign DNS/ANS resolution is proposed.  Real
identity authentication, credential verification, cryptographic handshakes,
blockchain immutability, post-quantum security, compliance certification,
universal injection prevention, and production isolation are not claimed.

The broader systems lesson is that trustworthy evidence must preserve negative
and incomplete outcomes.  A verifiable bundle is valuable precisely because it
can carry a failed policy decision without relabeling it as success.
